\section{Sublinear remainders in the lower norm comparison}

\begin{theorem}[Vanishing normalized logarithmic remainder]
  \label{thm:mean_tree_sublinear_remainder}
  \lean{Matrix.MeanTree.tendsto_rough_lower_remainder_div_nat}
  Fix real numbers $s,G,t,c,C_f$ and a real sequence $(\varepsilon_k)$
  such that $\varepsilon_k/k\to0$ as $k\to\infty$. Put
  \begin{align*}
    r_k&=\log(k+1)+c\log(k+2)+2t\varepsilon_k,\\
    E_k&=2sG\bigl(2|r_k|+C_f\log(k+2)+2\log2+\log(3/2)\bigr).
  \end{align*}
  Then $E_k/k\to0$. No sign condition on the five fixed coefficients
  or on $\varepsilon_k$ is required.

  This is the scalar limit for the explicit remainder in the rough
  lower norm estimate of \cite[Section~8]{OpenAI2026AreaLaw}.
\end{theorem}
\begin{proof}
  For every fixed real $a$, the limit
  $\log(k+a)/k\to0$ follows by substituting $x=k+a$ in
  $\log x/(x-a)\to0$. The cases $a=1,2$, together with the assumption
  on $\varepsilon_k$, give $r_k/k\to0$. Continuity of the absolute value
  gives $|r_k|/k\to0$. The remaining logarithmic and constant terms
  also vanish after division by $k$; multiplication by the fixed
  coefficient $2sG$ preserves the limit.
\end{proof}
